% page 70 of the facsimile (image p-069 of the scan)
% block 162.6 x 242.0 mm ; left margin 8.0 mm
\pgc{-12.700mm}{0.000mm}{
\l{}
\l{}
\l{\cel{8}62}
\l{}
\l{}
\l{\sou{benediktano}. Reguliero en la ordeno di santa Benediktus. -}
\l{\cel{2}DEFIRS.}
\l{}
\l{\sou{benefico}. (\sou{religio katolika)}. Ofico ekleziala qua implikis}
\l{\cel{3}la granto di revenuo.- DEFIRS.}
\l{}
\l{\sou{benigna}. Bonvoloza ad omna personi - indulgema ad omni. - DEFIS.}
\l{}
\l{\sou{benko}. Sidilo longa (ligna, kom exemplo) kun o sen dors-apogilo,}
\l{\cel{3}streta, sur qua plura personi povas sidar kune. - DEFIRS.}
\l{}
\l{\sou{benzino}.\cel{2}Hidrokarbido, liquido senkolora, extraktita de re-}
\l{\cel{3}zini e de gudro minkarbonala, qua uzesas por preparar granda}
\l{\cel{3}quanto de kompozaji aromatoza, por netigar la stofi, e c.}
\l{\cel{3}\sou{Simbolo kemiala}\cel{1}: C6 H6. - DEFIRS.}
\l{}
\l{\sou{benzoo}. Substanco rezinoza qua defluas, tra incizuro, de varie-}
\l{\cel{3}tato de L.\cel{1}\sou{styrax benzoin}.\cel{2}- DEFIRS.}
\l{}
\l{\sou{benzolo.}\cel{1}Kompozanto di la karbongudri qua distilo pasas sub}
\l{\cel{3}l7Oo C. e konsistas ye mixuro qua kontenas benzeno, toluene,}
\l{\cel{3}xilono. - DEFIRS.}
\l{}
\l{\sou{bequadro}. (\sou{muziko}) Signo analoga a la\cel{1}\sou{b}\cel{1}kun ventro quadrata,}
\l{\cel{3}uzata por nihiligar la bemoli o la diezi, e retroduktar}
\l{\cel{3}noto a lua tono naturala (de kande la diezi uzesis).-FIRS.}
\l{}
\l{\sou{bero}. (\sou{bot}.) Frukto tre pulpoza, di qua la semini situesas}
\l{\cel{3}meze di la pulpo. - DE.}
\l{}
\l{\sou{berberiso}. (\sou{bot.}) Planto lignoza, kovrita ye pikili, e qua}
\l{\cel{3}produktas grapi de beri reda tre acida, de qui onu facas}
\l{\cel{3}sorto di vino. - L.\cel{1}\sou{berberis vulgaris}.- DEFIRS.}
\l{}
\l{\sou{beredo}. Toko de lano, di qua la parto supra esas plata e cir-}
\l{\cel{3}kla. - DEFIR.}
\l{}
\l{\sou{berenjeno}. (\sou{bot.}) Varietato di solano, di qua la frukto (qua}
\l{\cel{3}havas la formo di granda bero longa) manjesas kom legumo.}
\l{\cel{3}L.\cel{1}\sou{solanum melongena}. - FS.}
\l{}
\l{\sou{bergamoto}. (\sou{bot.}) Varietato di piro tre parfumoza, qua quaze}
\l{\cel{3}liquideskas lor manjeso. L.\cel{1}\sou{citrus bergamia}. - DEFIRS.}
\l{}
\l{\sou{beriberio}. (\sou{Patol.}) Sorto di anemio, partikulara ye la habi-}
\l{\cel{3}tanti di insulo Ceylon. - F.}
\l{}
\l{\sou{berilo}. Varietato di smeraldo, diafana, maro-verda. - DEFIRS.}
\l{}
\l{\sou{berjero}. Larja fotelo, di qua la sido-plako esas garnisita}
\l{\cel{3}ye kuseno mola. - DFR.}
}
